\begin{itemframe}{مسئله استخدام}
\itm
فرض کنید شما می‌خواهید به کمک یک آژانس کاریابی دستیار جدیدی استخدام کنید.
آژانس کاریابی هر روز یک متقاضی را برای شما می‌فرستد و شما تصمیم می‌گیرید که آیا او را استخدام کنید یا خیر.
برای مصاحبه با هر متقاضی باید مبلغی به آژانس پرداخت کنید،
اما استخدام یک متقاضی هزینه‌ی بیشتری دارد، چرا که باید دستیار فعلی را اخراج کرده و هزینه‌ی قابل توجهی به آژانس بپردازید.
\itm
فرض کنید برای اینکه در هر زمان، بهترین فرد ممکن را داشته باشید، این روش را اتخاذ می‌کنید؛ گر هر متقاضی از دستیار فعلی بهتر بود، جایگزین دستیار فعلی شود.
\end{itemframe}


\begin{itemframe}{مسئله استخدام}
\itm
الگوریتم این رویکرد در شبه کد زیر قابل مشاهده است. متقاضیان از $1$ تا $n$ شماره‌گذاری شده‌اند.
\begin{algorithm}[H]\alglr
\caption{HIRE-ASSISTANT}
\begin{algorithmic}[1]
\Procedure{HIRE-ASSISTANT}{$n$}
  \State best $\gets 0$ \Comment{candidate 0 is a least-qualified dummy candidate}
  \For{$i \gets 1$ to $n$}
    \State interview candidate $i$
    \If{candidate $i$ is better than candidate $best$}
      \State $best \gets i$
      \State hire candidate $i$
    \EndIf
  \EndFor
\EndProcedure
\end{algorithmic}
\end{algorithm}
\itm
متقاضی شماره ۰ ساختگی است و از تمام متقاضیان دیگر صلاحیت کمتری دارد.
\end{itemframe}


\begin{itemframe}{مسئله استخدام}
\itm
مشخصاً این روش پر هزینه است. فرض کنید شما حاضر به پرداخت این هزینه هستند اما می‌‌خواهید تخمینی از این هزینه به دست آورید.
\itm
برخلاف تحلیل پیچیدگی زمانی ما تمرکز خود را نه بر زمان اجرای الگوریتم، بلکه بر هزینه‌هایی که برای مصاحبه و استخدام پرداخت می‌شود، قرار می‌دهیم.
\itm
در نگاه اول، تحلیل هزینه‌ی این الگوریتم ممکن است با تحلیل زمان اجرا متفاوت به نظر برسد.
با این حال، تکنیک‌های تحلیل یکسان هستند، چه در تحلیل هزینه و چه در تحلیل زمان اجرا. در هر دو حالت، ما تعداد دفعاتی که عملیات‌های پایه‌ای خاصی انجام می‌شوند را می‌شماریم.
\end{itemframe}


\begin{itemframe}{مسئله استخدام}
\itm
فرض کنید هزینه‌ی مصاحبه، $c_i$ و هزینه‌ی استخدام، $c_h$ باشد.
اگر $m$ تعداد افرادی باشد که استخدام می‌شوند، آنگاه هزینه‌ی کلی برابر است با:
$$
O(c_i n + c_h m)
$$
\itm
ما همیشه $n$ متقاضی را مصاحبه می‌کنیم پس همواره هزینه‌ی $c_i n$ را متحمل می‌شویم.
بنابراین، تمرکز خود را بر تحلیل هزینه‌ی استخدام یعنی $c_h m$ می‌گذاریم که به ترتیب متقاضیان بستگی دارد.


\end{itemframe}


\begin{itemframe-s}{مسئله استخدام}{تحلیل بدترین حالت}
\itm
در بدترین حالت، شما هر متقاضی‌ای که مصاحبه می‌کنید را استخدام می‌کنید. این حالت زمانی رخ می‌دهد که متقاضیان به ترتیب افزایشی از نظر صلاحیت ظاهر شوند؛ در این صورت، شما $n$ بار استخدام می‌کنید و هزینه‌ی کلی استخدام برابر است با:

$$
O(c_h n)
$$

\itm
البته، متقاضیان همیشه به ترتیب افزایشی نمی‌آیند. در واقع، شما هیچ اطلاعاتی درباره‌ی ترتیب آن‌ها ندارید و کنترلی نیز بر آن ندارید. بنابراین، طبیعی است که بپرسید در حالت معمول یا متوسط چه اتفاقی می‌افتد.
\end{itemframe-s}


\begin{itemframe-s}{مسئله استخدام}{تحلیل احتمالاتی}
\itm
تحلیل احتمالاتی
\fn{Probabilistic analysis}،
استفاده از احتمالات برای تحلیل مسائل است.
از تحلیل احتمالاتی برای تحلیل زمان اجرای الگوریتم و گاهی نیز برای تحلیل کمیت‌های دیگر، مانند هزینه‌ی استخدام در مسئله فعلی استفاده می‌شود.
\itm
برای انجام یک تحلیل احتمالاتی، باید دانشی درباره‌ی توزیع ورودی‌ها داشته باشیم یا فرضیاتی درباره‌ی آن‌ انجام دهیم.
سپس الگوریتم‌مان را تحلیل می‌کنیم و زمان اجرای میانگین آن را محاسبه می‌کنیم؛ یعنی به‌جای بررسی بدترین حالت یا بهترین حالت، میانگین زمان اجرا را روی همه ورودی‌های ممکن در نظر می‌گیریم. این میانگین با استفاده از محاسبه امید ریاضی به دست می‌آید.
\itm
به زمان اجرایی که از این طریق به دست می‌آید، زمان اجرای حالت میانگین
\fn{average-case running time}
 می‌گوییم.
\end{itemframe-s}


\begin{itemframe-s}{مسئله استخدام}{تحلیل احتمالاتی}
\itm
برای برخی مسائل، می‌توان فرض‌های معقولی درباره‌ی توزیع تمام ورودی‌های ممکن در نظر گرفت.
اما برای دیگر مسائل، نمی‌توان توزیع معقولی برای ورودی‌ها در نظر گرفت؛ در این موارد، تحلیل احتمالاتی ممکن نیست.
\itm
برای مسئله‌ی استخدام، می‌توان فرض کرد که متقاضیان به ترتیب تصادفی وارد می‌شوند. با استفاده از این دانش چطور می‌توانیم توزیع ورودی‌ها را مشخص کنیم؟

\itm
بیایید به هر متقاضی یک رتبه‌ی یکتا از $1$ تا $n$ بر اساس صلاحیت اختصاص دهیم. رتبه‌ی متقاضی $i$ را با $rank(i)$ نمایش می‌دهیم و فرض می‌کنیم که رتبهٔ بالاتر به معنای صلاحیت بیشتر است. (مثلاً اگر
$rank(3) = 1$
 باشد، یعنی متقاضی شماره‌ی ۳ بهترین است.)

\end{itemframe-s}


\begin{itemframe-s}{مسئله استخدام}{تحلیل احتمالاتی}
\itm
تعداد $n!$ حالت برای ترتیب رتبه‌ها وجود دارد.
این‌که متقاضیان به ترتیب تصادفی وارد می‌شوند، معادل آن است که بگوییم هر یک از این $n!$ حالت، با احتمال مساوی ظاهر می‌شوند.
 به‌عبارتی، رتبه‌ها یک جایگشت تصادفی یکنواخت
\fn{uniform random permutation}
 تشکیل می‌دهند.
\itm
به طور خلاصه؛ وقتی افراد تصادفی وارد می‌شوند، یعنی ترتیب رتبه‌ها تصادفی و یکنواخت است. یعنی همهٔ حالت‌های ممکن برای چیدن رتبه‌ها، با احتمال مساوی رخ می‌دهند.
\end{itemframe-s}


\begin{itemframe-s}{مسئله استخدام}{الگوریتم‌های تصادفی}
\itm
اما در بسیاری از موارد، اطلاعات اندکی درباره‌ی توزیع ورودی‌ها داریم. حتی اگر دانشی هم داشته باشیم، ممکن است مدل‌سازی محاسباتی آن امکان‌پذیر نباشد.
با این حال افزودن رفتار تصادفی به بخشی از الگوریتم می‌تواند کمک کننده باشد.
\itm
در مسئله‌ی استخدام، ممکن است به نظر برسد که متقاضیان به ترتیب تصادفی ارائه می‌شوند، اما هیچ راهی برای اطمینان از آن ندارید. بنابراین، برای توسعه‌ی یک الگوریتم تصادفی برای این مسئله، باید کنترل بیشتری بر ترتیب متقاضیان داشته باشید.
\itm
پس مدل را اندکی تغییر می‌دهیم: آژانس کاریابی فهرست تمام $n$ متقاضی را از قبل به شما می‌دهد. در هر روز، شما به صورت تصادفی انتخاب می‌کنید که با کدام متقاضی مصاحبه کنید. اگرچه درباره‌ی متقاضیان اطلاعاتی ندارید، اکنون به‌جای پذیرش ترتیبی که آژانس به شما می‌دهد و امید داشتن به تصادفی بودن آن، شما کنترل فرآیند را در دست دارید و ترتیبی تصادفی را اعمال می‌کنید.

\end{itemframe-s}


\begin{itemframe-s}{مسئله استخدام}{الگوریتم‌های تصادفی}
\itm
به‌طور کلی، یک الگوریتم را تصادفی
\fn{randomized}
می‌نامیم اگر رفتار آن تنها به ورودی وابسته نباشد، بلکه به مقادیری که یک تولیدکننده‌ی اعداد تصادفی
\fn{random-number generator}
نیز تولید می‌کند، بستگی داشته باشد.
\itm
فرض می‌کنیم که یک تابع تولیدکننده‌ی اعداد تصادفی به نام
\texttt{RANDOM}
در اختیار داریم. یک فراخوانی به صورت
$\texttt{RANDOM}(a, b)$
عددی صحیح بین $a$ و $b$ (شامل هر دو) بازمی‌گرداند، به‌طوری که هر عدد با احتمال مساوی ظاهر می‌شود.
\end{itemframe-s}


\begin{itemframe-s}{مسئله استخدام}{الگوریتم‌های تصادفی}
\itm
برای تحلیل زمانی الگوریتم‌های تصادفی، امید ریاضی زمان اجرا را بر روی توزیع مقادیر بازگردانده‌شده توسط تولیدکننده محاسبه می‌کنیم.
\itm
این الگوریتم‌ها را از الگوریتم‌هایی که ورودی آن‌ها تصادفی است، متمایز می‌کنیم؛ و زمان اجرای آن‌ها را «زمان اجرای مورد انتظار»
\fn{expected running time}
می‌نامیم.
\itm
به‌طور کلی، وقتی توزیع احتمالاتی در ورودی‌ وجود داشته باشد، در مورد «زمان اجرای حالت میانگین»،
و وقتی که خود الگوریتم تصمیمات تصادفی می‌گیرد، در مورد «زمان اجرای مورد انتظار» بحث می‌کنیم.
\end{itemframe-s}